\section{Escenarios de despliegue}

La validación del sistema se ha llevado a cabo en dos escenarios de despliegue que
cubren los casos de uso más representativos: un despliegue en dos PCs en red local y
un despliegue completamente contenerizado mediante Docker. Ambos escenarios utilizan
la misma configuración de voctocore y la misma interfaz gráfica, diferenciándose
únicamente en cómo se gestiona el entorno de ejecución.

\subsection{Escenario 1 --- Dos PCs en red local}

En este escenario, voctocore y las fuentes FFmpeg se ejecutan en un primer PC
(\textit{servidor/cerebro}) mientras que voctogui se conecta desde un segundo PC
mediante la red local. Este es el escenario que más se aproxima a una producción real
con cámaras físicas: el servidor de procesamiento puede situarse en un rack técnico
mientras el realizador trabaja desde otro equipo en la mesa de producción.

\textbf{Requisitos de instalación:} ambos equipos deben tener instalado Python 3,
las dependencias de GStreamer y sus plugins (\texttt{gstreamer1.0-plugins-base},
\texttt{-good}, \texttt{-bad}, \texttt{-ugly}, \texttt{libav}) y FFmpeg. La instalación
manual de estas dependencias puede ser tediosa, especialmente en distribuciones Linux
con versiones de GStreamer distintas a las esperadas.

Los scripts de arranque se encuentran en el directorio \texttt{2pc\_escenario2/}:
\begin{itemize}
  \item \texttt{arrancar\_voctocore\_pc1.sh}: inicializa el servidor en el PC de
        procesamiento, detectando la IP local y mostrándola para que el realizador
        la configure en su equipo.
  \item \texttt{arrancar\_voctogui\_pc2.sh}: lanza la interfaz gráfica en el PC del
        realizador, conectándose al servidor por su dirección IP.
\end{itemize}

La GUI recibe los flujos de previsualización en formato JPEG comprimido, minimizando
el ancho de banda requerido entre ambos equipos sobre una red local 1~GbE estándar.

\subsection{Escenario 2 --- Docker Compose}

El despliegue contenerizado elimina la necesidad de instalar manualmente Python ni
ninguna dependencia de GStreamer en el sistema anfitrión: todo el entorno de ejecución
se encuentra empaquetado dentro de la imagen Docker. El sistema completo se despliega
con un único comando desde el directorio \texttt{docker\_escenario3/}:

\begin{verbatim}
docker compose up -d
\end{verbatim}

Docker Compose levanta los diez servicios definidos en \texttt{docker-compose.yml} en
el orden correcto, gestionando automáticamente las dependencias de arranque mediante
los \textit{healthchecks} configurados. El realizador se conecta desde un segundo
equipo exactamente igual que en el Escenario 1.

Las ventajas de este escenario respecto al anterior son:
\begin{itemize}
  \item \textbf{Cero dependencias en el host:} basta con tener Docker instalado; no
        se requiere Python ni GStreamer.
  \item \textbf{Reproducibilidad total:} la imagen Docker garantiza que el entorno
        de ejecución es idéntico independientemente del sistema operativo anfitrión.
  \item \textbf{Gestión simplificada:} un único comando arranca y detiene todos los
        servicios, incluyendo voctocore, las fuentes, RabbitMQ y el servicio de
        telemetría.
\end{itemize}

% --- Figura sugerida ---
% \begin{figure}[H]
%   \centering
%   \includegraphics[width=0.9\textwidth]{figuras/diagramas_escenarios.png}
%   \caption{Diagramas de red de los dos escenarios de despliegue.}
%   \label{fig:diagramas_escenarios}
% \end{figure}
